\section{Partition functions as spatial moments}\label{sec:transfer}

We now construct a self-adjoint spatial operator whose moments are the
partition functions \eqref{eq:shift}. This identity lets even spatial
powers play the same role as positive Gibbs weights. The operator itself
need not be positive.

Let $\mathcal G\subset SO(3)$ be the group of proper signed permutation
matrices. It contains $D_2$ and $C$. For $g\in\mathcal G$, write
\[
 g e_\alpha=\epsilon_g(\alpha)e_{\sigma_g(\alpha)},
 \qquad \epsilon_g(\alpha)\in\{-1,1\}.
\]
Rotation covariance of the Cartesian matrices gives
\begin{equation}\label{eq:rotation-covariance}
 gG_\alpha g^{-1}
   =\epsilon_g(\alpha)G_{\sigma_g(\alpha)}.
\end{equation}

\begin{proposition}[Spatial transfer]\label{prop:transfer}
For each $b>0$ there are a Hilbert space $\mathcal K_b$, a compact
self-adjoint operator $X_b$, and a unitary representation
$g\mapsto U_g$ of $\mathcal G$ commuting with $X_b$, such that
\begin{equation}\label{eq:spatial-trace}
 Z_n(b,g)=\Tr(X_b^nU_g)
 \qquad(n\ge4,\ g\in\mathcal G).
\end{equation}
The operator $X_b$ and the representation commute with complex
conjugation. Moreover $X_b$ is Hilbert--Schmidt, so $X_b^k$ is trace class
for every integer $k\ge2$. If $\lambda_i(b)$ are its nonzero eigenvalues,
counted with multiplicity, then
\begin{equation}\label{eq:spatial-moments}
 \sum_i|\lambda_i(b)|^k<\infty\quad(k\ge2),\qquad
 Z_n(b)=\sum_i\lambda_i(b)^n\quad(n\ge4).
\end{equation}
\end{proposition}

\begin{proof}
\emph{Histories and the kernel.}
A bond history is a finite ordered list
\[
 x=((t_1,\alpha_1),\ldots,(t_k,\alpha_k)),
 \qquad 0<t_1<\cdots<t_k<b,\quad
 \alpha_j\in\{x,y,z\},
\]
including the empty list when $k=0$. Let $\Omega_b$ be the disjoint union
of these labeled simplexes. Give each simplex Lebesgue measure and the
empty history mass one. The resulting measure $\mu$ is finite:
\begin{equation}\label{eq:history-mass}
 \mu(\Omega_b)=\sum_{k=0}^\infty\frac{3^kb^k}{k!}=e^{3b}.
\end{equation}
Set $\mathcal K_b=L^2(\Omega_b,\mu;\C)$.

For two histories $x,y$, merge their labeled times in increasing order
and define $k_b(x,y)$ as the trace of the resulting product of the
matrices $G_\alpha$, with increasing time from left to right. A nonempty
time collision is a null event; set the kernel to zero there. The product
for two empty histories is the identity, so its trace is three.

The kernel is real and has absolute value at most three, since
$\norm{G_\alpha}=1$. Exchanging $x$ and $y$ leaves their merged
chronological list unchanged. Hence $k_b(x,y)=k_b(y,x)$; no reversal of
the matrix product occurs. Its integral operator $K_b$ is therefore
self-adjoint and Hilbert--Schmidt, with
\[
 \norm{K_b}_{\mathrm{HS}}^2\le9\mu(\Omega_b)^2=9e^{6b}.
\]
Define $X_b=e^{-ab}K_b$. Compact self-adjoint spectral theory gives real,
square-summable eigenvalues. Since they are bounded, all their absolute
moments of integer order at least two are summable. This proves the
operator assertions about $X_b$.

\emph{Rotations.}
The permutation $\sigma_g$ acts on the labels of a history, preserving
its times and the measure. Put
\[
 s_g(x)=\prod_{\alpha\text{ in }x}\epsilon_g(\alpha),\qquad
 (U_g f)(x)=s_g(\sigma_g^{-1}x)f(\sigma_g^{-1}x).
\]
The rule
$s_{gh}(x)=s_g(\sigma_hx)s_h(x)$ proves that these real unitaries form a
representation. Simultaneously rotating every matrix inside a site
trace, using \eqref{eq:rotation-covariance}, gives
\[
 k_b(\sigma_gx,\sigma_gy)=s_g(x)s_g(y)k_b(x,y).
\]
It follows that $U_g$ commutes with $K_b$ and $X_b$. In particular the
kernel of $K_bU_g$ is
\begin{equation}\label{eq:twisted-kernel}
 (K_bU_g)(x,y)=k_b(x,\sigma_gy)s_g(y).
\end{equation}

\emph{The cycle trace.}
For $n\ge2$, both $K_b^{n-1}$ and $K_bU_g$ are Hilbert--Schmidt.
The kernel of $K_b^{n-1}$ is the iterated kernel obtained by integration;
it is bounded because $k_b$ is bounded and $\mu$ is finite. The
Hilbert--Schmidt product trace identity
\[
 \Tr(AB)=\int a(x,y)b(y,x)\,d\mu(x)\,d\mu(y)
\]
and \eqref{eq:twisted-kernel} consequently give
\begin{equation}\label{eq:cycle-integral}
 \Tr(K_b^nU_g)=
 \int_{\Omega_b^n}
 \left(\prod_{j=0}^{n-2}k_b(x_j,x_{j+1})\right)
 k_b(x_{n-1},\sigma_gx_0)s_g(x_0)\,d\mu^n.
\end{equation}
The product trace identity follows from the Hilbert--Schmidt inner
product formula and Cauchy--Schwarz. All the iterated kernel integrals
here are absolutely integrable. Thus \eqref{eq:cycle-integral} does not
require a diagonal-kernel formula for $K_b$ itself.

\emph{Identification with the physical partition function.}
The ordered bond expansion is a standard representation of spin
partition functions~\cite[arXiv v1, Section~2.2, Equations~(2.10)--(2.12)]{AizenmanNachtergaele1994}.
We verify it here with the required seam and normalization.
Expand $\exp[-bH_n(g)]$ in its ordered-time series. By
\eqref{eq:cartesian-bond}, each event chooses a bond and one of three
labels, and inserts the corresponding two $G$ factors. At the seam the
site-$0$ factor is rotated as in \eqref{eq:rotation-covariance}.
There are $(3n)^k$ choices at total degree $k$; each tensor operator has
norm at most one. The sum of the absolute traced contributions is bounded
by
\begin{equation}\label{eq:history-absolute}
 3^n\sum_{k=0}^\infty\frac{(3nb)^k}{k!}=3^ne^{3nb}.
\end{equation}
We may therefore regroup all terms and integrals by bond. The
interleavings of the bond histories partition the global time simplex
up to null collisions. At each site, the tensor trace is the kernel on
the two incident histories, in their original chronological order.

For an explicit check of the seam, call its history $e_0$ and enumerate
the remaining bond histories $e_1,\ldots,e_{n-1}$ around the ring. The
physical integrand is
\[
 s_g(e_0)k_b(\sigma_ge_0,e_1)
 k_b(e_1,e_2)\cdots k_b(e_{n-1},e_0).
\]
In \eqref{eq:cycle-integral}, set
$x_0=e_0,x_1=e_{n-1},\ldots,x_{n-1}=e_1$ and use kernel symmetry.
This gives exactly the displayed integrand, with $g$ on site $0$.
The spatial enumeration has been reversed; the chronological products
have not. Thus $\Tr e^{-bH_n(g)}=\Tr(K_b^nU_g)$. Multiplication by
$e^{-abn}$ proves \eqref{eq:spatial-trace}. The ordinary trace formula
for a compact self-adjoint operator now gives
\eqref{eq:spatial-moments}.
\end{proof}

For an even $n\ge4$, Proposition~\ref{prop:transfer} supplies the
probabilities
\begin{equation}\label{eq:spatial-probabilities}
 p_i=\frac{|\lambda_i(b)|^n}{Z_n(b)},\qquad
 \sum_i p_i=1,\qquad
 \sum_i p_i^2=\frac{Z_{2n}(b)}{Z_n(b)^2}.
\end{equation}
Odd moments will also be useful in polynomial estimates, but evenness is
essential in \eqref{eq:spatial-probabilities}. These spatial eigenvalues
are distinct objects from the physical energies of $H_n$.

\subsection{Rotation sectors}

The small twisted partition functions separate the spatial eigenvalues
into a few lists. Their multiplicities will help isolate a dominant
weight; no restriction on the physical spectrum is imposed.

\begin{lemma}[Sector moments]\label{lem:sectors}
Fix $b>0$ and let $X_b,U_g$ be as in Proposition~\ref{prop:transfer}.
There are real eigenvalue lists, counted with multiplicity, whose
integer moments $I_k,O_k,N_k$ converge absolutely for $k\ge2$, with
\begin{align}
 Z_n(b)&=I_n+2O_n+3N_n,\label{eq:sector-total}\\
 Z_n(b,P)&=I_n+2O_n-N_n,\label{eq:sector-pi}\\
 Z_n(b,C)&=I_n-O_n\qquad(n\ge4).\label{eq:sector-cyclic}
\end{align}
The $I$ list is the restriction to the joint invariant subspace of
$D_2$ and $C$. The $O$ list occurs twice and the $N$ list three times
in the full eigenvalue list of $X_b$. All three moments are nonnegative
at even orders. If $J_n$ denotes the moment on the $D_2$-invariant
subspace, then
\[
 J_n=I_n+2O_n.
\]
\end{lemma}

\begin{proof}
The four characters $\chi$ of $D_2$ define orthogonal projections
\[
 \Pi_\chi=\frac14\sum_{g\in D_2}\chi(g)U_g.
\]
They commute with $X_b$ and sum to the identity. Conjugation by $U_C$
permutes the three nontrivial character spaces cyclically. The
restrictions of $X_b$ to these spaces are therefore unitarily equivalent;
denote their common moment by $N_k$.

In the trivial-character space, split further according to the
eigenvalues $1,\omega,\overline\omega$ of $U_C$, where
$\omega=e^{2\pi i/3}$. These are reducing subspaces for $X_b$. Complex
conjugation exchanges the last two spaces and preserves every real
eigenvalue of $X_b$, with multiplicity. Their common moment is $O_k$;
the moment on the first space is $I_k$. Absolute convergence follows
from \eqref{eq:spatial-moments}.

This orthogonal decomposition proves \eqref{eq:sector-total}. On the
three nontrivial character spaces, a fixed $U_P$ has signs $+1,-1,-1$,
while it is the identity on the trivial-character space. This gives
\eqref{eq:sector-pi}. The operator $U_C$ permutes the nontrivial spaces,
so they contribute zero to $\Tr(X_b^nU_C)$. On the remaining three
spaces its coefficients are $1,\omega,\overline\omega$, and
$\omega+\overline\omega=-1$. This proves
\eqref{eq:sector-cyclic}. The identity for $J_n$ follows from the
decomposition of the trivial-character space; even powers give nonnegativity.
\end{proof}

The invariant $I$ space need not be one-dimensional. The point of
Lemma~\ref{lem:sectors} is that every weight outside it occurs at least
twice in the full list. Later moment bounds will identify a unique
largest weight without discarding any sector.
